% =====================================================================
%  APERÇU DE RÉDACTION COMPLÈTE DU MÉMOIRE — FICHIER DE TRAVAIL
% =====================================================================
%  Statut : ce fichier N'EST PAS le rapport. Il vit à la racine du
%  projet, hors de rapport/, et ne doit jamais partir dans le rendu.
%
%  Usage : c'est un modèle de prose, chapitre par chapitre, dans l'ordre
%  du rapport. Là où tu avais déjà rédigé (intro, ch. 2, étape 1 du
%  ch. 3), ta prose est reprise telle quelle, seulement lissée. Le reste
%  est rédigé à partir de tes trames et de tes chiffres réels. Tu
%  restes maître : pioche, modifie, réécris dans les vrais fichiers.
%
%  Les tableaux et figures ne sont pas dupliqués ici : un repère
%  % [tableau/figure X ici] indique leur place. Les formules sont
%  recopiées quand la prose s'y appuie.
%
%  Chiffres : base 40 362 trades (11 stratégies, oracle inclus) pour les
%  ch. 1-3. ATTENTION ch. 4 : chiffres de l'ancienne base (38 035),
%  encore valides tant que les scripts ne sont pas relancés, mais à
%  régénérer SANS l'oracle avant le rendu (cf. TODO).
% =====================================================================


% =====================================================================
%  CHAPITRE 1 — INTRODUCTION
% =====================================================================

% ------------------ §1.1 But du mémoire (ta prose, lissée) -----------

Lors de mon stage en entreprise, j'ai eu pour consigne de créer un agent
d'intelligence artificielle capable de faire du trading, c'est-à-dire
d'acheter et de revendre des actions en cherchant à tirer profit de leurs
variations. J'ai alors décidé de donner à cet agent des capacités
d'analyse statistique. Après avoir construit le nécessaire pour qu'il
puisse passer des ordres et accéder aux données du marché, j'ai estimé
qu'il devait être capable d'évaluer les stratégies de trading qu'il
aurait établies — une stratégie étant une règle fixe qui décide quand
acheter un titre et quand le revendre.

Je présente dans ce rapport l'outil d'évaluation de stratégies que j'ai
donné à mon agent.

L'objectif de cet outil n'est pas de découvrir une stratégie rentable.
Il est de construire un enchaînement de méthodes capable de déterminer,
pour une stratégie quelconque qu'on lui soumet, si l'avantage qu'elle
affiche est réel ou s'il relève du hasard d'échantillonnage ou du biais
de marché. Autrement dit, la contribution revendiquée est
méthodologique~: la chaîne de tests, sa justification et la vérification
de ses conditions d'application.

Les dix stratégies d'analyse technique utilisées dans ce rapport ne
constituent donc pas l'objet d'étude, mais un banc d'essai. Elles
fournissent des jeux de données réels permettant de montrer que l'outil
discrimine correctement, qu'il sait rejeter ce qui ne tient pas et
retenir ce qui résiste à l'analyse. Les résultats obtenus sont présentés
à titre d'illustration du bon fonctionnement du protocole, et non comme
des recommandations d'investissement. Pour s'assurer que l'outil ne
rejette pas tout par excès de sévérité, on lui soumettra aussi un témoin
dont l'avantage est réel par construction.

% ------------------ §1.2 Les trois menaces (ta prose, lissée) --------

La première difficulté tient à la dérive du marché. Le premier réflexe
pour juger une stratégie est d'examiner si elle est rentable~: un gain
positif la ferait « marcher ». Ce raisonnement est pourtant trompeur,
car encore faut-il que ce gain dépasse ce qu'une entrée au hasard aurait
rapporté. Ici les stratégies sont appliquées aux titres du S\&P~500,
l'indice des cinq cents grandes capitalisations américaines, dont la
tendance est haussière sur la période étudiée. Cette dérive haussière
pose problème, car une stratégie peut réaliser des gains sans avoir
capté le moindre signal. C'est pourquoi on ne mesurera pas les gains,
mais l'avantage qu'apporte chaque stratégie par rapport au hasard.

La deuxième difficulté provient du nombre de tests. Plus on teste de
stratégies, plus le risque d'obtenir au moins un faux positif augmente.
Ce problème sera traité au chapitre~\ref{chap:juger} à l'aide de la
correction de Bonferroni.

La troisième, sans doute la plus délicate, est la dépendance des
données. De nombreuses méthodes statistiques exigent l'indépendance des
observations~; or ici les trades sont dépendants sous plusieurs aspects,
comme la période, le titre ou encore le secteur. Cette difficulté sera
mesurée au chapitre~\ref{chap:donnees}, puis neutralisée au
chapitre~\ref{chap:juger}.

On est forcé de prendre en compte ces trois difficultés, sous peine d'un
verdict trop optimiste~: c'est leur traitement conjoint qui fait de ce
protocole un véritable outil de validation.

% ------------------ §1.3 Démarche du mémoire (rédigé, était en trame)

L'étude suit le fil d'une enquête unique, du diagnostic au verdict. Le
chapitre~\ref{chap:donnees} présente d'abord la matière première, la
base des quarante mille trades et la construction de l'avantage mesuré,
puis dresse l'état des lieux de ce qui menace les tests, en mesurant les
trois canaux par lesquels les observations se ressemblent. Le
chapitre~\ref{chap:juger} constitue le c\oe ur du travail. On y juge les
stratégies, du test de référence jusqu'aux corrections de la dépendance,
en resserrant la prudence à chaque étape, avant de relire le verdict à
la lumière du comportement du marché lui-même. Le
chapitre~\ref{chap:contexte} change ensuite de question~: il ne s'agit
plus de valider, mais de comparer les stratégies entre elles et de
décrire les situations dans lesquelles elles se déclenchent. Le
chapitre~\ref{chap:protocole} récapitule enfin l'outil sous la forme
d'un protocole réutilisable, en assume les limites et trace les étapes
suivantes, avant la conclusion. Le fil conducteur va ainsi du plus
fragile, la donnée et ses dépendances, au plus solide, le verdict, puis
l'outil qui permettra d'en rendre d'autres.


% =====================================================================
%  CHAPITRE 2 — LES DONNÉES ET LEUR STRUCTURE DE DÉPENDANCE
% =====================================================================

% ------------------ §2.1 La base des trades (ta prose, lissée) -------

Tout au long de ce rapport, on appliquera les méthodes à notre base de
données. Celle-ci a été obtenue en appliquant différentes stratégies à
des données antérieures. Cette méthode se nomme le \emph{backtesting}~:
elle consiste à rejouer le passé pour vérifier si une stratégie a repéré
un signal de prédiction du marché.

Ici, nous appliquons dix stratégies techniques classiques ainsi qu'une
stratégie témoin, appelée oracle. La stratégie oracle utilise les
données futures pour effectuer des trades uniquement gagnants, afin de
vérifier que notre protocole laisse passer une stratégie dont le signal
fonctionne réellement. Les dix autres stratégies, elles, effectuent
leurs trades à une date fixée en utilisant uniquement les données
précédant cette date.

% [tableau tab:strategies ici — les onze stratégies]

J'ai alors choisi d'appliquer ces stratégies aux cinq cent une grandes
capitalisations américaines présentes dans le S\&P~500. Cet indice
constitue un ensemble sur lequel on peut investir avec peu de risque~:
même si le cours de chacune de ces actions peut chuter individuellement,
l'indice lui-même reste relativement stable et suit une tendance
haussière.

Pour qu'une stratégie soit validée, il faut qu'elle obtienne de
meilleurs résultats qu'une stratégie jouant au hasard, et que cela ne
soit pas dû à notre échantillonnage. Or la tendance est ici haussière,
ce qui nous ramène à la première difficulté de l'introduction~: comment
évaluer une stratégie lorsque le marché monte de lui-même~?

Pour répondre à cette question, on simule le hasard. À chaque trade
réalisé, on associe deux cents trades fictifs de même exposition, sur le
même titre, avec la même durée et le même sens, mais dont la date
d'entrée est tirée au sort. La moyenne de leurs rendements définit le
gain du hasard pour ce trade. On définit alors, pour chaque trade, son
\edgevar{} comme le rendement effectivement obtenu par la stratégie,
éventuellement négatif, diminué de ce gain du hasard~:
\[
  \edgevar = \text{gain de la stratégie} - \text{gain du hasard}.
\]
Un \edgevar{} positif signifie que la stratégie a fait mieux qu'une
entrée au hasard de même exposition~: son \emph{timing} apporte une
information que la seule dérive du marché n'explique pas.

On obtient ainsi une base contenant, pour chaque stratégie, entre 978 et
8\,513 trades effectués entre janvier~2010 et juillet~2026, pour un
total de 40\,362 trades. Chaque ligne correspond à une position ouverte
puis fermée par une stratégie sur un titre. On y associe le titre
concerné, désigné par son \emph{ticker}, c'est-à-dire son code en
bourse, la stratégie employée, le secteur d'activité, les dates d'entrée
et de sortie, la durée de détention, le régime de marché et le contexte
à l'entrée, le rendement obtenu et l'avantage \edgevar{} sur un tirage
aléatoire équivalent.

% [tableau tab:dico-trades ici — variables + exemple]

% ------------------ §2.2 Les rendements sectoriels (ta prose, lissée)

Pour étudier la dépendance des données, on a besoin de données par
secteur. Le S\&P~500 compte onze secteurs d'activité. Sur chacun d'eux,
on établit une courbe du rendement journalier, obtenu en moyennant
chaque jour les rendements des titres du secteur. Ces données seront
utilisées d'une part au canal~2 du
paragraphe~\ref{sec:dep-canal2}, pour mettre en évidence le facteur
marché, c'est-à-dire la part des mouvements commune à tous les
secteurs, et d'autre part dans la conclusion du
chapitre~\ref{chap:juger}, pour examiner si le marché est prévisible à
partir de son propre passé.

% ------------------ §2.3 chapeau (ta prose, lissée) ------------------

Les tests que nous utiliserons au chapitre~\ref{chap:juger} ne
fonctionnent que si les données sont indépendantes. Ici, les données
sont dépendantes par le titre, par la date et par la durée. Cette
situation porte un nom en théorie des sondages~: notre base n'est pas un
tirage au hasard d'individus, mais un tirage par grappes, où chaque
grappe est un titre. Le réflexe du sondeur impose alors de décrire ce
plan d'échantillonnage avant d'estimer quoi que ce soit. C'est pourquoi
on mesure les dépendances maintenant, avant les tests.

% ------------------ Canal 1 — la dépendance de titre (ta prose, lissée)

Deux trades d'un même titre sont liés par la gestion de la société ainsi
que par son image à l'extérieur. Ils subissent les mêmes chocs propres à
la société, ses résultats, ses décisions de direction, ses litiges,
indépendamment du moment où ils sont ouverts. On appellera cet effet la
dépendance de titre~: nos données constituent alors un tirage par
grappes, où chaque grappe est un titre.

On définit alors l'estimateur suivant~:
\[
  \hat\rho \;=\; \frac{\mathrm{MSB}-\mathrm{MSW}}
                      {\mathrm{MSB} + (n_0 - 1)\,\mathrm{MSW}},
\]
où MSB (\emph{Mean Square Between}) est la variation entre les titres,
donc inter-grappes, c'est-à-dire la moyenne des écarts d'\edgevar{} d'un
titre à l'autre, et MSW (\emph{Mean Square Within}) la variation à
l'intérieur d'un titre, donc intra-grappe, c'est-à-dire la moyenne des
écarts d'\edgevar{} entre les trades d'un même titre. Deux cas se
présentent. Si les trades d'un même titre ne se ressemblent pas plus que
ceux de titres différents, la variation est la même à l'intérieur des
titres et entre eux, et l'on obtient MSB~$\approx$~MSW. S'ils se
ressemblent, c'est la dépendance de titre, la variation à l'intérieur
d'un titre s'effondre tandis que la variation entre titres subsiste, et
l'on obtient MSB~$>$~MSW.

On pourrait se contenter du quotient de MSB par MSW, qui vaudrait un en
cas d'indépendance. Mais un écart à un, qu'il soit de un dixième, de
trois ou de cent, ne se rattache à aucune échelle, et l'on ne saurait
dire s'il traduit une dépendance forte ou faible. Pour obtenir une
réponse du type « treize pour cent du résultat est dû à une dépendance
de grappe », on préfère la formule ci-dessus, qui ramène la mesure à une
valeur comprise entre zéro et un.

Le terme $n_0$ correspond à la taille moyenne des grappes, c'est-à-dire
au nombre de trades par titre, corrigée de l'inégalité des tailles. Il
pondère la variation intra-titre au dénominateur~: plus les grappes
contiennent de trades, plus $n_0$ est grand, et plus il faut de
contraste entre les titres pour conclure à une dépendance.

Le tableau~\ref{tab:dep-mesures} donne un $\hat\rho$ quasi nul pour la
plupart des stratégies, mais atteignant $0{,}134$ pour
\texttt{rsi\_strict}. La dépendance par le titre est donc négligeable
pour toutes les stratégies, excepté \texttt{rsi\_strict}. Un $\hat\rho$
faible ne garantit toutefois pas l'indépendance des trades~: il n'écarte
que ce premier canal, les deux suivants restant à examiner.

% ------------------ Canal 2 — la dépendance de marché (ta prose, lissée)

La deuxième dépendance que l'on va mesurer est la dépendance du moment.
Des trades ouverts au même moment peuvent être liés parce qu'ils
subissent la même journée de bourse~: à une même période, différentes
actions changent de valeur conjointement, et les trades effectués sont
alors liés par le marché.

Cette dépendance ne s'observe pas directement sur les trades, mais sur
le marché lui-même. Si les secteurs bougent ensemble, alors des trades
ouverts au même moment, quel que soit leur secteur, subissent la même
poussée. Pour la mesurer, on utilise donc les données introduites au
paragraphe~\ref{sec:donnees-secteurs}, c'est-à-dire les rendements
quotidiens par secteur. On calcule les corrélations des secteurs deux à
deux, puis on en fait la moyenne, et l'on obtient une corrélation de
$0{,}70$. Ce chiffre élevé signifie que les secteurs sont largement
corrélés. Lorsqu'un secteur bouge, le marché a tendance à bouger dans le
même sens, et inversement, lorsque le marché suit une tendance, la
plupart des secteurs la suivent aussi.

On utilise ensuite la méthode de l'ACP sur les onze secteurs, de façon à
mettre en évidence les plus grandes corrélations
(figure~\ref{fig:acp-marche}). Cette méthode calcule les axes qui
captent la plus grande variance. En l'appliquant ici, on observe un
premier axe qui capte 73~\% de la variance totale, le long duquel tous
les secteurs progressent ensemble~: c'est le mouvement global du marché.
Vient ensuite un deuxième axe, nettement moins important, qui sépare les
secteurs cycliques, réagissant fortement au marché, des secteurs
défensifs, qui ne bougent guère quelle que soit la conjoncture. On peut
donner en exemple de secteur cyclique la technologie, qui dépend du
pouvoir d'achat, on achète des biens technologiques quand on est aisé,
moins en temps de crise, et comme secteur défensif les services aux
collectivités, électricité, eau et gaz, dont la consommation reste
constante, que le marché soit haussier ou en crise.

% [figure fig:acp-marche ici]

Sur l'histogramme de gauche, la première barre écrase toutes les
autres~: elle capte à elle seule 73~\% de la variance, contre moins de
8~\% pour la deuxième. Sur le nuage de droite, tous les secteurs se
situent du même côté de l'axe horizontal, ce qui confirme qu'ils
progressent ensemble, tandis que le second axe les répartit entre
défensifs et cycliques.

Cette cause commune n'est donc pas une hypothèse mais une mesure~: 73~\%
des mouvements sont portés par une force unique. Deux trades ouverts au
même moment subissent largement le même sort, ce qui les rend dépendants
du marché.

% ------------------ Canal 3 — la dépendance temporelle (ta prose, lissée)

Dans la suite du rapport, nous aurons besoin d'agréger les trades par
mois. Or les trades durent~: en moyenne 132~jours, avec une médiane de
65~jours, et plus de 75~\% d'entre eux dépassent un mois. Un même trade,
encore ouvert, est donc compté dans plusieurs mois consécutifs. Deux
mois voisins partagent ainsi des trades communs et se ressemblent. Même
agrégées par mois, les observations ne sont donc pas indépendantes, et
c'est cette ressemblance d'un mois à l'autre qu'il faut mesurer.

Pour cela, on utilise l'autocorrélation d'ordre un, notée
$\hat\rho(1)$~: on décale la série mensuelle d'un mois et on corrèle la
série avec elle-même ainsi décalée. Elle mesure à quel point un mois
ressemble au précédent. On obtient les valeurs suivantes.

% [tableau tab:dep-mesures ici — rho intra-titre + rho(1) mensuel]

Le $\hat\rho(1)$ atteint $0{,}56$ pour \texttt{ma\_crossover}, sa valeur
la plus forte, tandis qu'il reste faible pour \texttt{rsi\_strict}
($0{,}03$). La dépendance temporelle est a priori forte pour les
stratégies à trades longs, faible pour celles à trades courts.
Attention~: ces valeurs élevées ne traduisent pas une prévisibilité du
marché, mais le chevauchement mécanique des trades qui durent.

% ------------------ Conclusion du §2.3 (ta prose, lissée) ------------

Les trois canaux de dépendance sont désormais mesurés. Aucun n'est
systématiquement élevé. La dépendance de titre reste souvent proche de
zéro, et même l'autocorrélation temporelle s'efface pour les stratégies
à trades courts. Mais chacun se manifeste fortement pour certaines
stratégies, la dépendance de titre pour \texttt{rsi\_strict}, la
dépendance de marché pour toutes, la dépendance temporelle pour
\texttt{ma\_crossover}, et le chapitre~\ref{chap:juger} montrera qu'ils
suffisent à renverser des verdicts. On peut maintenant juger les
stratégies en sachant exactement ce qui menace les tests.


% =====================================================================
%  CHAPITRE 3 — JUGER LES STRATÉGIES : L'ESCALADE DE PRUDENCE
% =====================================================================

% ------------------ Chapeau du chapitre ------------------------------

Ce chapitre traite la question suivante~: parmi onze stratégies
d'analyse technique, lesquelles produisent un avantage \edgevar{}
significativement positif~? On applique d'abord le test de référence
sous ses hypothèses standard, puis on met à l'épreuve, une à une, ses
conditions d'application, en s'appuyant sur la structure de dépendance
mesurée au chapitre~\ref{chap:donnees}.

% ------------------ Étape 1 — la question (ta prose) -----------------

La première question posée à l'outil est donc~: l'avantage moyen affiché
par une stratégie est-il réel, ou n'est-il qu'un accident
d'échantillonnage~? On la formalise en confrontant la moyenne de
l'\edgevar{} à la valeur zéro, ce qui conduit au test de Student
présenté ci-dessous.

% ------------------ Étape 1 — Student (ta prose, lissée) -------------

Lorsque l'on obtient un edge moyen positif, on cherche à savoir si la
stratégie constitue réellement une prédiction, ou si ce résultat est dû
à l'aléa de l'échantillon. Il faut ici distinguer deux quantités~:
l'edge moyen théorique \(\mu_{\edgevar}\), c'est-à-dire la vraie valeur,
inconnue, qui gouverne la stratégie, et sa moyenne observée \(\bar e\),
calculée sur notre échantillon de trades. On se demande alors si, à
partir de la valeur observée \(\bar e\), on peut conclure que la vraie
valeur \(\mu_{\edgevar}\) est positive.

Le test de Student à un échantillon répond à cette question en
confrontant \(\mu_{\edgevar}\) à une valeur de référence fixée à zéro,
zéro correspondant à l'absence d'avantage sur le hasard. Sa mise en
\oe uvre suppose deux conditions, à savoir que l'edge vérifie un critère
de normalité et que les trades soient indépendants les uns des autres.
La première sera assouplie par le théorème central limite, la seconde,
plus délicate, sera examinée en section suivante. On oppose ainsi deux
hypothèses~:
\[
  \begin{cases}
    \Hzero : \mu_{\edgevar} = 0 & \text{(le vrai edge est nul~: pas mieux que le hasard),}\\[2pt]
    \Hun : \mu_{\edgevar} > 0 & \text{(le vrai edge est positif~: avantage réel).}
  \end{cases}
\]
Le test est unilatéral, car seul un avantage strictement positif nous
intéresse.

Pour trancher, on définit \(t\) comme le nombre d'erreurs-types
nécessaires pour faire passer l'edge moyen de zéro à celui observé dans
notre échantillon~:
\[
  t = \frac{\bar e}{s/\sqrt{n}} \sim t_{n-1} \quad\text{sous } \Hzero .
\]
Sous \Hzero{}, \(t\) suit une loi de Student à \(n-1\) degrés de
liberté. On évalue alors la p-value, c'est-à-dire la probabilité
d'obtenir un \(t\) au moins aussi grand que celui observé si \Hzero{}
était vraie. Si cette p-value est très faible, le hasard n'est plus une
explication crédible de l'edge moyen positif observé, et l'on conclut
que l'edge théorique est positif~: la stratégie fait mieux que le
hasard.

On a dit que le test de Student demandait la normalité de l'edge. On
vérifie cette condition d'abord avec le test de Shapiro--Wilk, puis à
l'aide du théorème central limite. Le test de Shapiro--Wilk a en effet
peu de chances de valider l'hypothèse~: la taille de l'échantillon
déstabilise ce test, alors qu'elle renforce au contraire le théorème
central limite.

% ------------------ Étape 1 — Shapiro (ta prose, lissée) -------------

Le test de Shapiro--Wilk est un test de normalité qui compare deux
mesures de dispersion, d'une part les écarts à la moyenne, d'autre part
l'ordre des valeurs. Ces deux mesures ne coïncident que si l'hypothèse
de normalité est vraie. On oppose ainsi deux hypothèses~:
\[
  \begin{cases}
    \Hzero \;(W \simeq 1) : \text{l'edge suit une loi normale,}\\[2pt]
    \Hun \;(W \ll 1) : \text{l'edge ne suit pas une loi normale.}
  \end{cases}
\]
On notera que, contrairement au test de Student, c'est ici \Hzero{} qui
correspond à la propriété recherchée~: une p-value faible conduira donc
à rejeter la normalité. On pose alors la statistique
\(W = b^2/\mathrm{SCE}\) comme rapport de ces deux mesures. Lorsque
\(W\) est proche de un, l'échantillon est compatible avec une loi
normale. Pour trancher, on procède comme pour le test de Student~: on
calcule la probabilité d'obtenir notre \(W\) si la loi était normale, et
si cette probabilité est trop faible, on rejette l'hypothèse de
normalité.

Appliqué à nos données, le test rejette la normalité pour les dix
stratégies, les valeurs de \(W\) et les p-values étant détaillées en
annexe~\ref{tab:ann-shapiro}. La condition du test de Student n'est
donc, en apparence, pas remplie.

% ------------------ Étape 1 — TCL (ta prose, lissée) -----------------

Le rejet de Shapiro--Wilk porte sur la non-normalité de l'edge
individuel. Mais comme le test de Student ne dépend que de la moyenne de
l'edge, on peut se rabattre sur le théorème central limite, qui garantit
la quasi-normalité de la moyenne sur les grands échantillons, quelle que
soit la loi de départ~:
\[
  \frac{\bar e - \mu_{\edgevar}}{\sigma/\sqrt{n}}
  \;\xrightarrow[n \to \infty]{}\; \mathcal{N}(0,1) .
\]
Cette justification suppose un effectif suffisant, et un outil générique
devrait écarter les stratégies trop peu fournies. Ici, \(n\) est compris
entre 978 et 8\,513 selon la stratégie, ce qui est largement suffisant.
La condition réellement utile au test de Student est donc satisfaite,
malgré le rejet de Shapiro--Wilk. On remarque de plus que la quantité
ci-dessus est exactement celle du test de Student~: sous \Hzero{}, en
remplaçant l'écart-type théorique \(\sigma\), inconnu, par son
estimation \(s\), on retrouve la statistique
\(t = \bar e/(s/\sqrt{n})\). Le théorème central limite justifie ainsi
directement la forme et la validité du test.

% ------------------ Étape 1 — Bonferroni (ta prose, lissée) ----------

Pour appliquer le test de Student, il reste à fixer un seuil pour la
p-value. Pour un test isolé, on prend habituellement \(0{,}05\)~: on
estime que s'il n'y avait que cinq chances sur cent d'obtenir un tel
résultat sous \Hzero{}, il est peu vraisemblable d'être tombé dessus par
hasard, et l'on accepte alors le risque de rejeter \Hzero{}.

Seulement, nous effectuons ici onze tests. Plaçons-nous dans le cas le
plus défavorable, où aucune stratégie n'a d'avantage réel. Pour un test
isolé, la probabilité de ne pas conclure à tort à un avantage est de
\(0{,}95\). En supposant les tests indépendants, la probabilité de ne
commettre aucune erreur sur les onze tests tombe à
\(0{,}95^{11} \approx 0{,}57\). Le risque de commettre au moins un faux
positif sur l'ensemble s'élève donc à \(1 - 0{,}95^{11} \approx
0{,}43\), ce qui est inacceptable.

Pour éviter cela, on raisonne en sens inverse~: on fixe le risque global
à \(0{,}05\) et l'on en déduit le seuil de chaque test. On s'appuie sur
l'inégalité de Boole, selon laquelle la probabilité qu'au moins un
événement se produise est inférieure ou égale à la somme de leurs
probabilités individuelles. Pour garantir un risque global inférieur à
\(0{,}05\), il suffit donc que la somme des seuils individuels vaille
\(0{,}05\). En les prenant tous égaux, on obtient la correction de
Bonferroni~:
\[
  \alpha' = \frac{\alpha}{m} = \frac{0{,}05}{11} \approx 0{,}0045 .
\]

% [encadré conditions ici — ajouter le renvoi : « l'indépendance des
%  trades est mise à l'épreuve en section suivante »]

% ------------------ Étape 1 — Résultats (ta prose, lissée) -----------

% [tableau tab:b1-etape1 ici]

Outre l'oracle témoin, qui bat le hasard de très loin comme attendu,
seules \texttt{rsi\_classic} et \texttt{rsi\_strict} survivent à la
correction de Bonferroni.

Il est utile de rappeler ce que ce test affirme au juste. Il répond
uniquement à la question de savoir si l'edge moyen positif est dû au
seul échantillonnage, auquel cas on conserve \Hzero{} et la stratégie
n'est pas déclarée gagnante. Dans le cas inverse, on pourra affirmer que
la stratégie a capturé un signal, mais pas qu'elle est gagnante, et cela
pour deux raisons. D'une part, un edge positif ne signifie pas des gains
positifs, mais simplement que l'on perd moins que le hasard. D'autre
part, comme \(n\) est grand, la loi normale de la moyenne se resserre,
et l'on fait ressortir comme significatives des stratégies dont la
moyenne d'edge n'est que très faiblement positive, que les frais
suffiraient à rendre perdantes.

% [encadré limites ici — significativité vs ampleur]

Ce verdict reste toutefois suspendu à une condition que nous n'avons pas
encore vérifiée, l'indépendance des trades. C'est l'objet de la section
suivante.

% ------------------ Étape 1 bis — la question (version fluide) -------

On a appliqué le test de Student en supposant les données indépendantes.
Or, comme nous l'avons mesuré au chapitre~\ref{chap:donnees}, nos
données sont dépendantes par le titre, par le marché et d'un mois à
l'autre. Ces dépendances empruntent en réalité deux chemins seulement,
que nous appellerons les deux portes. La première est la porte de
l'action, celle des trades ouverts sur un même titre. La seconde est la
porte de la période, celle des trades ouverts à une même époque, qu'ils
subissent ensemble le mouvement du marché ou qu'ils se prolongent sur
les mêmes mois. Ce sont les deux seules voies par lesquelles deux trades
peuvent partager une cause commune.

Le problème est que le test de Student traite les trades d'une stratégie
comme autant d'informations distinctes, alors que ces informations se
répètent en partie. L'erreur-type s'en trouve sous-estimée. Comme cette
erreur-type sert à mesurer l'écart entre l'avantage observé et la valeur
de référence, un dénominateur trop petit gonfle la statistique de
Student et rend la p-value trop optimiste. Le risque de conclure à tort
qu'une stratégie bat le hasard augmente d'autant.

Nous cherchons donc à savoir si le verdict positif obtenu pour
\texttt{rsi\_classic} et \texttt{rsi\_strict} résiste lorsque l'on
abandonne l'hypothèse d'indépendance. Pour cela, nous corrigeons d'abord
l'erreur-type au moyen de l'effet de plan, puis nous vérifions ce
résultat par un rééchantillonnage qui ne suppose aucune forme
particulière des données. Nous fermons ensuite chaque porte en
regroupant les trades pour que chaque titre, puis chaque période, ne
compte qu'une fois. Nous vérifions enfin que la porte de la période
tient elle-même, en traitant la suite des avantages mensuels comme une
série temporelle.

% ------------------ Étape 1 bis — méthode 1 : l'effet de plan --------

Nos trades ne forment pas un échantillon tiré au hasard parmi des
observations indépendantes. Ils proviennent d'un ensemble de titres,
chacun livrant plusieurs trades. La théorie des sondages appelle cela un
tirage par grappes, où chaque grappe est ici un titre. Le
chapitre~\ref{chap:donnees} a déjà mesuré, avec la corrélation
intra-titre $\hat\rho$, à quel point les trades d'une même grappe se
ressemblent. Il reste à en tenir compte dans le test.

L'effet de plan, introduit par Kish, chiffre exactement ce que cette
mise en grappes coûte à la précision. Il se calcule ainsi~:
\[
  \mathrm{DEFF} = 1 + (\bar m - 1)\,\hat\rho,
  \qquad n_{\mathrm{eff}} = \frac{n}{\mathrm{DEFF}},
  \qquad \widehat{\mathrm{se}}_{\mathrm{corr}}
        = \frac{s}{\sqrt{n}}\sqrt{\mathrm{DEFF}},
\]
où $\bar m$ désigne le nombre moyen de trades par titre. La variance
réelle de la moyenne vaut $\mathrm{DEFF}$ fois celle que l'on
obtiendrait sous indépendance. Tout se passe donc comme si l'on ne
disposait que de $n_{\mathrm{eff}}$ trades réellement informatifs au
lieu de $n$. Lorsque la corrélation est nulle, l'effet de plan vaut un
et rien ne change. Lorsqu'elle augmente, l'erreur-type est gonflée et le
test devient plus exigeant. On reprend alors le test de Student de
l'étape~1 avec cette erreur-type corrigée, en conservant le même seuil
de Bonferroni.

% ------------------ Étape 1 bis — méthode 2 : le bootstrap -----------

Le test précédent s'appuie encore sur la loi de Student, donc sur une
forme particulière des données. Pour nous en affranchir, nous confirmons
le résultat par un rééchantillonnage, la méthode du bootstrap proposée
par Efron~\cite{efron1979}. Le principe consiste à retirer au sort, un
grand nombre de fois, un nouvel échantillon à partir de celui que l'on
possède. Ici nous tirons des titres entiers, avec remise, et jamais des
trades isolés. La structure de dépendance interne de chaque titre est
ainsi transportée telle quelle dans chaque tirage. Après quatre mille
répétitions, la proportion de tirages où l'avantage disparaît fournit
une p-value qui ne suppose ni normalité, ni indépendance entre les
trades d'un même titre. La seule hypothèse qui subsiste est
l'indépendance entre titres différents.

% ------------------ Étape 1 bis — méthode 3 : les deux portes --------

Corriger l'erreur-type ne suffit pas. L'effet de plan conserve en effet
la pondération du test initial, dans lequel un titre porteur de deux
cents trades pèse deux cents fois plus qu'un titre n'en portant qu'un
seul. Un avantage concentré sur quelques titres très actifs peut ainsi
passer le test alors qu'il ne se retrouve pas ailleurs. Pour écarter ce
défaut, nous changeons de logique et regroupons les trades avant de
tester, de sorte que chaque unité ne compte qu'une fois.

La première porte, celle de l'action, consiste à calculer d'abord
l'avantage moyen de chaque titre, puis à appliquer le test de Student à
ces quelque cinq cents moyennes. La seconde porte, celle de la période,
calcule de la même manière l'avantage moyen de chaque mois d'entrée,
puis teste ces quelque cent quatre-vingt-dix moyennes. Dans les deux
cas, la dépendance à l'intérieur d'une grappe ne peut plus rien fausser,
puisqu'un titre, ou un mois, ne pèse qu'une seule voix. Le théorème
central limite s'applique alors à ces moyennes de grappes, dont le
nombre reste élevé. Pour être validée, une stratégie doit franchir les
deux portes, chacune au seuil de Bonferroni.

% ------------------ Étape 1 bis — méthode 4 : la série mensuelle -----

La porte de la période repose à son tour sur une hypothèse, celle de
mois indépendants. Or nous avons vu au chapitre~\ref{chap:donnees} que
les trades se prolongent souvent au-delà d'un mois, si bien que deux
mois voisins partagent des positions encore ouvertes et se ressemblent.
La suite des avantages mensuels est donc autocorrélée, et le test
précédent pourrait rester trop optimiste. Pour lever ce dernier doute,
nous traitons cette suite comme une série temporelle, en suivant le
déroulement classique de ce type d'analyse.

Avant tout, la série doit être stationnaire, c'est-à-dire garder les
mêmes caractéristiques dans le temps. Nous le vérifions par deux tests
dont les hypothèses de départ sont opposées, celui de Dickey et Fuller
et celui dit KPSS. Lorsque les deux concordent, la conclusion ne dépend
pas du test choisi. Nous examinons ensuite l'autocorrélation de la série
pour déterminer combien de mois passés influencent le mois courant, ce
qui fixe l'ordre d'un modèle autorégressif que nous estimons.

Ce modèle fournit alors une variance de la moyenne qui tient compte de
la persistance, appelée variance de long terme~:
\[
  \gamma_{\mathrm{lr}}
     = \frac{\sigma^2_\eta}{\bigl(1-\sum_{i=1}^{p}\Phi_i\bigr)^{2}},
  \qquad
  \widehat{\mathrm{Var}}(\bar x) \approx \frac{\gamma_{\mathrm{lr}}}{T},
  \qquad
  \mathrm{DEFF}_t = \frac{\gamma_{\mathrm{lr}}}{\gamma(0)} .
\]
Cette variance additionne toutes les autocovariances de la série, là où
la formule usuelle en $s^2/T$ suppose au contraire des mois sans lien.
Plus la persistance est forte, plus la correction est importante. Nous
reprenons enfin le test de Student avec cette variance corrigée, puis
nous vérifions, par le test de Box et Pierce sur les résidus du modèle,
que celui-ci a bien capté toute l'autocorrélation. Si les résidus se
comportent comme un bruit sans mémoire, la correction est justifiée.

% [encadré conditions ici — indépendance entre grappes, TCL sur les
%  moyennes, stationnarité DF+KPSS, Box-Pierce]

% ------------------ Étape 1 bis — Résultats : lecture 1b -------------

% [tableau tab:b1-dep-deff ici]

La corrélation intra-titre étant quasi nulle pour \texttt{rsi\_classic}
et modérée pour \texttt{rsi\_strict}, la correction par l'effet de plan
ne change presque rien~: les deux stratégies conservent leur verdict
positif, et le bootstrap par grappes confirme ces p-values sans recourir
à la loi de Student. À ce stade, la porte de l'action semble donc
franchie. Ce test reste cependant pondéré par le nombre de trades, et ne
dit pas si l'avantage est partagé par l'ensemble des titres ou concentré
sur quelques titres très actifs. C'est ce que le test suivant va
trancher.

% ------------------ Étape 1 bis — Résultats : lecture 1c (le cœur) ---

% [tableau tab:b1-dep-portes ici]

La lecture de ce tableau constitue le résultat central du mémoire.
L'oracle, notre témoin, franchit les deux portes de très loin, avec des
p-values proches de zéro. L'outil sait donc dire oui lorsque l'avantage
est réel~: son éventuelle sévérité ne condamne pas d'avance toutes les
stratégies.

Le sort de \texttt{rsi\_classic} est tout autre. La stratégie tombe aux
deux portes, avec des p-values de $0{,}059$ et $0{,}039$, loin du seuil
de $0{,}0045$. Le \(t\) de l'étape~1 valait pourtant $5{,}6$, mais ses
7\,117 trades ne portaient pas 7\,117 informations indépendantes. Dès
que chaque titre ne vote qu'une fois, l'avantage disparaît~: il était
concentré sur les titres et les mois les plus tradés. Le verdict positif
de l'étape~1 était donc un faux positif, que l'outil vient de démasquer
en vérifiant ses propres conditions.

Le cas de \texttt{rsi\_strict} est plus subtil. La stratégie franchit
brillamment la porte de l'action, avec une p-value de
$3{,}9\times10^{-4}$~: son avantage est partagé par de nombreux titres.
Mais elle échoue à la porte de la période, avec une p-value de
$0{,}26$~: présente sur 138~mois seulement, elle concentre son avantage
sur quelques épisodes. Faute de régularité temporelle démontrable, son
signal n'est pas exploitable tel quel.

Au bilan, en dehors du témoin, plus aucune stratégie ne franchit les
deux portes. L'outil discrimine~: il valide l'avantage construit pour
être réel et rejette ceux qui ne tiennent que par la dépendance des
données.

% ------------------ Étape 1 bis — Résultats : lecture 1d -------------

% [tableau tab:b1-dep-mensuel ici]

Les moyennes mensuelles se révèlent bel et bien autocorrélées, jusqu'à
$0{,}56$ pour \texttt{ma\_crossover}~: l'hypothèse de mois indépendants,
sur laquelle reposait la porte de la période, était elle-même optimiste,
la variance étant sous-estimée d'un facteur atteignant $6{,}4$. Toutes
les séries sont stationnaires au vu des tests de Dickey--Fuller et KPSS,
et les modèles autorégressifs sont validés par le test de Box--Pierce
sur leurs résidus~: la correction est donc fondée. Après correction,
aucun verdict ne bascule. \texttt{rsi\_classic} termine à $0{,}08$, loin
du seuil. La porte de la période est fermée proprement. Signalons par
honnêteté que les mois sans trade sont traités comme consécutifs, avec
une couverture de 71~\% pour \texttt{rsi\_strict} et supérieure à 88~\%
ailleurs, limite assumée de cette construction.

% [encadré limites ici — Kish approché, perte de puissance assumée,
%  mois sans trade, ordre AR plafonné]

% ------------------ Étape 1 bis — positionnement littérature ---------

Les outils de cette section ont été construits à partir des seuls
concepts du cours, la théorie des sondages, l'analyse de la variance et
les séries temporelles. Il se trouve qu'ils rejoignent les standards de
l'économétrie financière, signe que la démarche converge vers la
pratique établie. L'agrégation par période, qui calcule une statistique
par mois avant de mener l'inférence sur la suite de ces statistiques,
est exactement la démarche proposée par Fama et
MacBeth~\cite{famamacbeth1973}, devenue un standard du domaine. La
variance de long terme estimée au travers du modèle autorégressif
poursuit le même objectif que les erreurs-types robustes à
l'autocorrélation de Newey et West~\cite{neweywest1987}, dont elle est
une version paramétrique, validée ici par le test de Box--Pierce. Enfin,
le contrôle des faux positifs sur un très grand univers de stratégies
relèverait du Reality Check de White~\cite{white2000}~; avec onze
stratégies déclarées à l'avance, la correction de Bonferroni suffit et
demeure le choix le plus prudent.

Le verdict au standard corrigé est donc le suivant~: aucune des dix
stratégies réelles ne démontre d'avantage. Le verdict de l'étape~1 ne
tenait que par une condition non vérifiée. C'est le résultat central de
ce mémoire~: l'outil de validation a démasqué son propre faux positif en
testant ses conditions d'application. Un verdict négatif rendu
proprement est la preuve que l'outil fonctionne.

% ------------------ Lecture du verdict (le marché lui-même) ----------

Le verdict est tombé, aucune stratégie ne démontre d'avantage. Avant de
conclure, on peut se demander si ce résultat était prévisible, en
interrogeant le marché lui-même avec les outils des séries temporelles.
Le chapitre~\ref{chap:donnees} a montré que la dépendance simultanée est
massive, le facteur marché portant 73~\% des mouvements. Reste à savoir
s'il existe de la dépendance retardée, la seule qu'une stratégie puisse
exploiter, puisqu'il faut un délai pour passer un ordre.

On vérifie d'abord la stationnarité, condition de toute modélisation,
par le test de Dickey--Fuller augmenté, appliqué au prix puis au
rendement de l'indice de marché. Le prix ne rejette pas la racine
unitaire, avec une p-value de $0{,}999$~: il se comporte comme une
marche aléatoire, dont le niveau de demain est celui d'aujourd'hui plus
un choc imprévisible. Le rendement, lui, est nettement stationnaire,
avec une p-value de l'ordre de $10^{-26}$. On modélise donc le
rendement, jamais le prix. Son autocorrélation est quasi nulle à tous
les retards examinés, entre $-0{,}08$ et $+0{,}09$ sur les dix premiers
jours. Un modèle ARIMA~\cite{boxjenkins1970} ajusté sur ce rendement ne
trouve d'ailleurs rien à saisir, ses deux termes, autorégressif et
moyenne mobile, se compensant presque exactement.

Le rendement du marché est ainsi proche d'un bruit sans mémoire~: son
propre passé ne prédit guère le lendemain. C'est ce que la littérature
appelle l'efficience faible du marché. La boucle est alors bouclée. Le
marché regorge de dépendance simultanée, celle qui fausse les tests et
qu'aucune stratégie ne peut exploiter faute de délai, et manque
précisément de dépendance retardée, la seule qui serait exploitable.
L'échec des dix stratégies n'est pas un accident~: c'est la signature de
cette asymétrie.

% [encadré conditions ici — rendements stationnaires, jamais les prix]

% ------------------ Synthèse du chapitre 3 ---------------------------

Retraçons le chemin parcouru. Le test de référence déclarait deux
stratégies gagnantes sur dix. La mise à l'épreuve de ses conditions
d'application, par l'effet de plan, le bootstrap par grappes,
l'agrégation en deux portes et la correction de la série mensuelle, a
renversé ces deux verdicts. \texttt{rsi\_classic} était un faux positif
né de la répétition d'une même information sur des milliers de trades.
\texttt{rsi\_strict} possède un avantage réel à l'échelle des titres,
mais trop épisodique dans le temps pour être retenu. La lecture du
marché lui-même rend ce résultat cohérent~: un marché proche de
l'efficience faible n'offre presque pas de dépendance retardée à
exploiter. Le message central est ailleurs que dans les stratégies. Le
protocole a fait exactement son travail, il a démasqué son propre faux
positif en vérifiant ses conditions, tout en validant sans hésiter le
témoin dont l'avantage était réel. Un verdict négatif rendu proprement,
c'est l'outil qui fonctionne.


% =====================================================================
%  CHAPITRE 4 — LE CONTEXTE DES TRADES : COMPARER ET CARACTÉRISER
% =====================================================================
%  NB : chiffres de l'ancienne base — à régénérer SANS l'oracle (TODO).

% ------------------ Chapeau du chapitre ------------------------------

Le chapitre précédent a jugé chaque stratégie isolément, contre le
hasard. Celui-ci change de question. Les stratégies diffèrent-elles
entre elles~? Le succès dépend-il du contexte, du régime de marché, du
secteur, de la situation d'entrée~? Les conclusions de ce chapitre sont
comparatives et descriptives~: on n'y valide aucune stratégie, le
verdict appartenant au chapitre~\ref{chap:juger}.

% [encadré : exclusion de l'oracle, déjà rédigé dans le fichier]

% ------------------ ANOVA (liaison + résultats) ----------------------

Pour comparer les stratégies entre elles, on analyse la variance de
l'\edgevar{}~: si les moyennes des dix stratégies étaient égales, la
variation entre stratégies ne dépasserait pas la variation interne à
chacune. La statistique \(F\) rapporte ces deux variations~:
\[
  F = \frac{\mathrm{SCE}_{\text{inter}}/(k-1)}
           {\mathrm{SCE}_{\text{intra}}/(n-k)} \sim F_{k-1,\,n-k}.
\]
L'analyse à un facteur rejette massivement l'égalité des moyennes, avec
\(F\approx 47\) et une p-value de l'ordre de $10^{-85}$. Le test de
Tukey, fondé sur la loi du range studentisé et qui corrige de lui-même
la multiplicité des paires, identifie dix-huit paires significatives sur
quarante-cinq, \texttt{rsi\_strict} se détachant du lot. Une analyse à
deux facteurs croisant la stratégie et le régime de marché révèle enfin
une interaction très significative, de l'ordre de $5\times10^{-69}$~: la
valeur d'une stratégie dépend du régime dans lequel elle opère. Ces
p-values, calculées au niveau des trades, s'entendent sous la réserve
d'indépendance du chapitre~\ref{chap:juger}~; les \(F\) colossaux
restent interprétables comme ordres de grandeur.

% ------------------ Khi-deux (liaison + résultats, ta base) ----------

On se demande ensuite si le simple fait de gagner dépend du contexte. La
variable \texttt{win}, qui vaut un pour un trade gagnant, est croisée
avec le régime, le secteur et la stratégie dans des tableaux de
contingence, sur lesquels on applique le test du khi-deux
d'indépendance~:
\[
  \chi^2 = \sum_{i,j}\frac{(N_{ij}-E_{ij})^2}{E_{ij}},
  \qquad E_{ij}=\frac{N_{i\cdot}N_{\cdot j}}{N}.
\]

% [tableau tab:b1-etape3 ici]

Le taux de gain dépend fortement de la stratégie, faiblement du régime,
et pas du secteur. À très grand \(n\), le khi-deux devient toutefois
significatif même pour des liens négligeables~: la significativité ne
dit rien de la force du lien, et ces contrastes se lisent comme des
descriptions, non comme des validations.

% ------------------ ACP du contexte d'entrée (lecture rédigée) -------

Pour cartographier les situations d'entrée, on réalise une ACP normée
sur les quatre variables quantitatives de contexte, la volatilité, le
RSI, la distance à la moyenne mobile et la durée de détention. Les deux
premiers axes résument environ deux tiers de la variance, le premier
étant porté par le RSI et la distance à la moyenne mobile, autrement dit
par l'état de survente du titre à l'entrée, le second par la durée de
détention et la volatilité. La question intéressante est de savoir si
les trades gagnants occupent une région particulière de ce plan. La
réponse est négative~: les positions moyennes des gagnants et des
perdants sont presque confondues, avec un écart standardisé inférieur à
$0{,}2$. Le contexte d'entrée, à lui seul, ne sépare pas les gagnants
des perdants, résultat instructif qui rejoint le verdict du
chapitre~\ref{chap:juger}.

% [figure etape4_acp ici]

% ------------------ AFC et ACM (lectures rédigées) -------------------

L'analyse factorielle des correspondances croise ensuite la stratégie
avec la tranche de rendement obtenue, perte, gain modéré ou gain fort.
Le premier axe, qui résume les trois quarts de l'inertie, oppose les
stratégies à dominante perdante, comme \texttt{hs\_classic} ou
\texttt{sr\_breakout}, aux stratégies dont les trades finissent plus
souvent en gain, \texttt{rsi\_classic} et surtout \texttt{rsi\_strict},
cette dernière se projetant nettement du côté des gains forts. L'analyse
des correspondances multiples, menée sur la stratégie, le régime, le
secteur et l'issue du trade, ne résume que faiblement l'information,
comme il est habituel avec de nombreuses modalités, mais confirme la
même géographie~: les modalités de secteur se massent au centre du plan,
signe qu'elles ne structurent rien, tandis que les stratégies et le
régime s'écartent, \texttt{rsi\_strict} voisinant avec le régime
baissier, dont elle exploite les excès.

% ------------------ Synthèse du chapitre 4 ---------------------------

Les stratégies diffèrent donc fortement entre elles et leur valeur
dépend du régime de marché, tandis que le secteur ne joue aucun rôle
décelable. Mais aucun de ces contrastes ne remet en cause le verdict du
chapitre~\ref{chap:juger}~: différer de son voisin n'est pas battre le
hasard. Ces cartes du contexte préparent en revanche la perspective d'un
avantage conditionnel, esquissée au chapitre~\ref{chap:protocole}.


% =====================================================================
%  CHAPITRE 5 — LE PROTOCOLE, SES LIMITES, SES SUITES
% =====================================================================

% ------------------ Le protocole récapitulé --------------------------

Ce que le stage livre n'est pas une stratégie gagnante, mais une
procédure de jugement, réutilisable telle quelle sur tout signal futur
de l'agent, qu'il soit technique, informationnel ou relationnel. Le
tableau~\ref{tab:protocole} récapitule cette procédure en six étages.
Chaque étage répare une menace précise identifiée au fil du mémoire, de
la dérive du marché à l'autocorrélation résiduelle, et le verdict n'est
prononcé que si tous les étages concordent. C'est cette exigence
conjointe qui fait la valeur de l'ensemble~: la défaillance d'une seule
condition suffit, comme on l'a vu, à fabriquer un faux positif.

% [tableau tab:protocole ici — les six étages]

% ------------------ Limites assumées ---------------------------------

Aucune des limites qui suivent n'invalide le verdict rendu. Chacune
trace en revanche la frontière de ce que le protocole peut affirmer.
Les résultats sont établis sur l'échantillon qui a servi aux calculs, et
une validation hors échantillon reste nécessaire avant tout usage réel.
Le jugement porte sur dix stratégies déclarées à l'avance, et un univers
beaucoup plus vaste exigerait des outils dédiés au risque de fouille de
données, comme le Reality Check de White~\cite{white2000}. Les frais et
conditions d'exécution sont modélisés simplement. Enfin, la puissance
élevée des grands échantillons impose de distinguer soigneusement la
significativité statistique de la taille de l'effet. On peut ajouter que
la base ne contient que les entreprises présentes dans l'indice sur la
période~: les sociétés disparues en sont absentes, biais dit du
survivant, atténué ici par la construction appariée de l'\edgevar{}, qui
compare chaque trade au hasard sur le même titre.

% ------------------ Perspectives -------------------------------------

La première suite naturelle est la validation temporelle. Elle consiste
à figer les règles sur une période d'apprentissage puis à juger sur la
période suivante, en avançant pas à pas, ce que la littérature appelle
le walk-forward. La validation croisée classique est ici interdite,
puisqu'elle mélangerait le futur et le passé.

La deuxième est l'avantage conditionnel. Les cartes du
chapitre~\ref{chap:contexte} suggèrent que le succès dépend du régime de
marché~; modéliser la probabilité de gain en fonction du contexte, par
exemple par une régression logistique, en serait le prolongement direct.
Cette voie a été écartée du présent mémoire par rigueur~: sur des
dizaines de milliers de trades dépendants, les p-values des coefficients
seraient trop optimistes, exactement le piège du
chapitre~\ref{chap:juger}. Les outils adaptés, modèles à effets mixtes
ou équations d'estimation généralisées, dépassent le cadre de ce
travail et en constituent la prochaine étape. Une version outillée du
protocole devrait aussi intégrer un garde-fou d'effectif, écartant en
amont les stratégies trop peu fournies pour que le théorème central
limite s'applique.

La troisième concerne les autres signaux de l'agent. Celui-ci collecte
également des signaux d'information, contrats publics, transactions de
responsables politiques, annonces réglementaires, et des relations
d'avance ou de retard entre titres. Des travaux exploratoires ont été
engagés sur ces deux fronts. Le protocole du chapitre~\ref{chap:juger}
est précisément l'outil qui permettra de les juger, hors du périmètre de
ce mémoire.


% =====================================================================
%  CHAPITRE 6 — CONCLUSION
% =====================================================================

L'objectif fixé en introduction n'était pas de découvrir une stratégie
gagnante, mais de construire un protocole capable de juger n'importe
laquelle. Ce protocole est construit, éprouvé, et il a rendu son
verdict~: aucune des dix stratégies d'analyse technique ne démontre
d'avantage réel une fois ses conditions d'application vérifiées.

Le résultat central mérite d'être rappelé. Le test standard déclarait
deux stratégies gagnantes. La vérification de la seule condition
d'indépendance a renversé les deux verdicts, l'un se révélant un faux
positif né de la répétition d'une même information sur des milliers de
trades, l'autre un avantage réel mais trop épisodique pour être retenu.
Dans le même temps, le témoin construit pour être gagnant a été validé
sans hésitation à chaque étage. Un outil de validation se juge à sa
capacité d'attraper ses propres erreurs tout en reconnaissant le vrai~:
la démonstration est faite.

Une idée traverse l'ensemble de ce travail. La dépendance des données y
a joué les deux rôles. Simultanée, portée par le facteur marché qui
explique près des trois quarts des mouvements, elle était la nuisance
qui faussait les tests, et il a fallu la mesurer puis la neutraliser
pour juger proprement. Retardée, elle aurait été la seule ressource
qu'une stratégie puisse exploiter, et le marché n'en contient presque
pas, conformément à l'efficience faible. Neutraliser l'une, chercher
l'autre~: c'est la même statistique, vue des deux côtés du même objet.

Le protocole est désormais prêt à juger les prochains signaux de
l'agent, en commençant par leur validation temporelle.
